\documentclass[10pt,a4paper]{amsart}
\usepackage[margin=19mm]{geometry}
\usepackage{amsmath,amssymb,mathtools}
\usepackage[scr=rsfs,cal=euler]{mathalfa}
\usepackage{xfrac}
\usepackage[unicode,hidelinks,bookmarks=false]{hyperref}
\newcommand{\ie}{i.e.\ }
\newcommand{\cf}{cf.\ }
\newcommand{\resp}{respectively\ }
\newcommand{\Subsection}{\subsection}
\providecommand{\nobreakdash}{\nobreak\hbox{-}}
\setlength{\emergencystretch}{2em}
\multlinegap0pt
\usepackage{amssymb}
\usepackage{mathtools}
\usepackage{bm}
\usepackage[shortlabels]{enumitem}
\newtheorem{lemma}{Lemma}
\newtheorem{proposition}{Proposition}
\usepackage{cleveref}
\crefname{lemma}{Lemma}{Lemmas}
\crefname{proposition}{Proposition}{Propositions}
\newcommand{\R}{\mathbb R}
\title{\textbf{Vorticity-Response Blowup in Axisymmetric Models}}
\author{Rishad Shahmurov}
\date{Source: arXiv:2604.09949v2 (CC BY 4.0)}
\begin{document}
\maketitle

\noindent\emph{Source and licence.} \textbf{Vorticity-Response Blowup in Axisymmetric Models}. Version arXiv:2604.09949v2 (CC BY 4.0), distributed by arXiv under the Creative Commons Attribution 4.0 International licence, http://creativecommons.org/licenses/by/4.0/, as recorded in the arXiv licence metadata for this exact version and shown on the abstract page https://arxiv.org/abs/2604.09949v2. Full licence notice: \texttt{licenses/arxiv-2604.09949-CC-BY-4.0.txt}.\par\smallskip

\noindent\textit{Editorial scope.} Counted unit: the article's proposition prop:continue (source0.tex lines 832--839). The article's complete original proof of it is delivered with it (source0.tex lines 841--880, one \texttt{proof} environment of 40 lines). No mathematics is written, repaired or re-derived: the statement and the proof are the article's own lines, in the article's own notation, and no retained line is substituted or re-flowed.  Retained uncounted as the source-local context the proof needs: lines 328--334; lines 348--363; lines 371--374; lines 528--531.  Nothing was omitted from any retained span, so the package carries no omission record.  No retained line contains a citation, so the wrapper carries no bibliography and no bibliography entry.  No retained line contains a figure environment, an image-inclusion command or drawing code, so no image is delivered and the package declares no asset dependency.  Editorial formatting only, in addition: an AMS \texttt{amsart} preamble that restates the article's own declarations and macros used by the retained lines, namely the environments \texttt{lemma}, \texttt{proposition} on separate counters, together with the standard packages those lines need.  The retained environments print in this document as Lemma 1, Proposition 1 in order of appearance, whereas the article prints the same items with the numbering of its own full document.  The complete native source of the article as distributed in the arXiv e-print for this exact version is bundled unchanged as \texttt{sources/198236/source0.tex}.
\begin{equation}\label{eq:twochannel}
\begin{aligned}
 (a_+)_t+V(a_+)_x&=a_+^2+\nu(a_+)_{xx},\\
 (a_-)_t+V(a_-)_x&=a_-^2+\nu(a_-)_{xx},\\
 V_x&=a_++a_-,\qquad V(0,t)=0.
\end{aligned}
\end{equation}

\begin{lemma}\label{lem:cone}
Assume that $a_{+,0},a_{-,0}\in C_c^\infty(\R)$ are even, nonnegative, and nonincreasing on $(0,\infty)$.  Then, on the maximal classical lifespan, each $a_\pm$ remains even and nonnegative and satisfies
\begin{equation}\label{eq:monotone}
 (a_\pm)_x(x,t)\le0\qquad(x>0).
\end{equation}
Consequently
\begin{equation}\label{eq:Vsign}
 V(x,t)\ge0\qquad(x>0)
\end{equation}
and
\begin{equation}\label{eq:supersolution}
 (a_\pm)_t-\nu(a_\pm)_{xx}
 =a_\pm^2-V(a_\pm)_x
 \ge a_\pm^2.
\end{equation}
\end{lemma}

\begin{equation}\label{eq:hpm}
 (h_\pm)_t+V(h_\pm)_x
 =\nu(h_\pm)_{xx}+(a_\pm-a_\mp)h_\pm.
\end{equation}

\begin{equation}\label{eq:mass-id}
 \frac{d}{dt}\int a_j
 =2\int a_j^2+\int a_ja_k.
\end{equation}

\begin{proposition}\label{prop:continue}
Let $\nu>0$ and let $(a_+,a_-)$ be a smooth solution of \eqref{eq:twochannel} on $[0,T)$ satisfying the cone conclusions of \cref{lem:cone}.  If
\begin{equation}\label{eq:Astar}
 \sup_{0\le t<T}
 \bigl(\|a_+(t)\|_\infty+\|a_-(t)\|_\infty\bigr)<\infty,
\end{equation}
then the solution extends smoothly beyond $T$.
\end{proposition}

\begin{proof}
Let the bound in \eqref{eq:Astar} be $A$.  We first control the masses
\[
 M_j(t)=\int_\R a_j(x,t)\,dx.
\]
Because the channels are nonnegative, \eqref{eq:mass-id} gives
\[
 M_j'(t)
 =2\int a_j^2+\int a_ja_k
 \le (2\|a_j\|_\infty+\|a_k\|_\infty)M_j(t)
 \le 3A M_j(t).
\]
Hence
\begin{equation}\label{eq:massbound}
 M_j(t)\le M_j(0)e^{3At}.
\end{equation}
Since $V(0,t)=0$ and $V_x=a_++a_-$,
\begin{equation}\label{eq:Vbound}
 \|V(t)\|_\infty
 \le M_+(t)+M_-(t),
\end{equation}
so the drift stays bounded on $[0,T)$.

Next set $h_j=(a_j)_x$.  Equation \eqref{eq:hpm} is linear in $h_j$ once $a_\pm$ and $V$ are fixed.  The maximum principle applied to $h_j$ and $-h_j$ gives
\begin{equation}\label{eq:hbound}
 \|h_j(t)\|_\infty
 \le \|h_j(0)\|_\infty e^{2At}.
\end{equation}
We now differentiate repeatedly.  Let $D=\partial_x$ and $w_{j,n}=D^na_j$.  Applying $D^n$ to the $j$th equation gives
\begin{equation}\label{eq:higher}
 (w_{j,n})_t+V(w_{j,n})_x
 =\nu(w_{j,n})_{xx}
 +(2a_j-nV_x)w_{j,n}+P_{j,n},
\end{equation}
where $P_{j,n}$ is a finite sum of products of spatial derivatives of $a_+$ and $a_-$ of order strictly less than $n$.  This form follows because the only term containing $D^{n+1}a_j$ is $V D^{n+1}a_j$, and the only additional order-$n$ contribution from differentiating $V(a_j)_x$ is $nV_xD^na_j$.

Assume inductively that all derivatives of both channels through order $n-1$ are bounded on $[0,T)\times\R$.  Then $P_{j,n}$ is bounded, while the coefficient $2a_j-nV_x$ is bounded by a constant depending only on $n$ and $A$.  The parabolic maximum principle applied to \eqref{eq:higher}, followed by Gronwall, yields a uniform bound for $D^na_j$ on $[0,T)$.  Starting from \eqref{eq:hbound}, induction controls every spatial derivative.  The equations then control time derivatives as well.

These bounds are precisely the a priori estimates required by the standard local parabolic existence theorem to restart the solution at times tending to $T$.  Therefore the solution extends beyond $T$.
\end{proof}

\end{document}
